\documentclass[11pt]{article}
\usepackage[T1]{fontenc}
\usepackage{lmodern}
\usepackage[margin=1in]{geometry}
\usepackage{amsmath,amssymb,amsthm,booktabs,array}
\usepackage[hidelinks]{hyperref}
\newtheorem{theorem}{Theorem}
\newtheorem{lemma}[theorem]{Lemma}
\newtheorem{corollary}[theorem]{Corollary}
\newcommand{\one}{\mathbf{1}}
\title{Hardin and Barker's strip-coloring formulas:\\
an extension of Britnell and Wildon's enumeration}
\author{Haobo Ma\\\small ChronoAI Pte Ltd, Singapore
\and Wenlin Zhang\\\small National University of Singapore\\\small ChronoAI Pte Ltd, Singapore}
\date{16 November 2026}
\begin{document}
\maketitle
\begin{abstract}
We extend Britnell and Wildon's marked-partition enumeration (2015,
Proposition 8.1) to prove Colin Barker's empirical generating function in
the FORMULA entry of 11 October 2018 in OEIS A233254, revision 8, for Hardin's
normalized six-column arrays at every positive height, by three-color transfer.
An ordered-graph normalization formula and an eight-state matrix prove its
associated recurrence. The method also proves the neighboring formulas and
all empirical row and column recurrences listed in Hardin's table A233256,
revision 6, through seven rows or seven columns.
The larger fixed-height cases use explicit finite integer-matrix identities.
\end{abstract}

\section{The six-column formula}
Let $a_n$ count $n\times6$ arrays $c_{r,j}\in\{0,1,2,3,4,5\}$, with
rows $1\le r\le n$ and columns $1\le j\le6$, subject to the following rules.
Horizontally adjacent entries differ and do not sum to $5$; the same rule
holds between $c_{r,j+1}$ and $c_{r+1,j}$ for $r<n$, $j<6$.
There are no vertical or main-diagonal restrictions.
In row-major order the first entry is $0$; each occurrence of $2$, $3$, or $4$
has an earlier $1$, and each $3$ has an earlier $2$.
Set $a_0=0$.
These are Hardin's arrays in OEIS A233254 \cite{six}; Barker gave the
empirical generating function there. The following theorem supplies a
finite-state proof, with every matrix entry specified below.
Transfer matrices for chromatic polynomials are standard; see
Salas and Sokal \cite[Section 3.1]{salas}.
Here the normalization formula and eight boundary shapes give the stated
coefficients without a six-color first-occurrence state space.
Britnell and Wildon \cite[Proposition 8.1 and Section 6.3]{britnell}
previously proved the marked-partition interpretation of the one-column
sequence cross-referenced by Hardin's table. We give an ordered-graph version
of that enumeration and use it for the bounded-width and bounded-height
formulas. Their one-column result is credited separately below; it is not
a prior proof of Barker's six-column generating function.

\begin{theorem}\label{thm:main}
As an identity of formal power series over $\mathbb Q$,
\begin{equation}\label{eq:gf}
\sum_{n\ge1}a_nx^n=
\frac{8x(17-7488x+974848x^2-40108032x^3)}
{(1-64x)(1-1600x+368640x^2-21233664x^3)}.
\end{equation}
Consequently, for $n\ge5$,
\begin{equation}\label{eq:rec}
a_n=1664a_{n-1}-471040a_{n-2}+44826624a_{n-3}-1358954496a_{n-4},
\end{equation}
with $a_1=136$, $a_2=166400$, $a_3=220626944$, and $a_4=294517735424$.
\end{theorem}

\subsection{Removing the color normalization}
For a finite simple graph $G$ with an ordered, nonempty vertex set of size $N$,
write $P_G(q)$ for its number of proper colorings using $q$ available colors.
Impose the preceding first-occurrence rules on six-colorings of its vertices,
and exclude equal or complementary colors on every edge.

\begin{lemma}\label{lem:normalization}
The number $A(G)$ of these normalized colorings is
\begin{equation}\label{eq:general}
A(G)=\frac{2^N}{48}\bigl(P_G(3)+3P_G(2)+15P_G(1)\bigr).
\end{equation}
If $G$ is connected and bipartite with at least one edge, this becomes
$A(G)=2^N(P_G(3)+6)/48$.
\end{lemma}
\begin{proof}
Replace each color by its complementary pair $\{0,5\}$, $\{1,4\}$, or
$\{2,3\}$. Every color restriction is exactly the assertion that the pair
labels at the endpoints differ. Let $b_k$ count unordered partitions of the
vertices into $k$ nonempty independent sets. Once such a partition is fixed,
order its sets by their first vertices. Their first colors are respectively
$0,1,2$; the other vertices each have two independent color choices in their
assigned pair. Thus
\[
A(G)=2^{N-1}b_1+2^{N-2}b_2+2^{N-3}b_3.
\]
This is a bijection: the first-occurrence rules select exactly these pair
names and first colors, including when fewer than three pairs occur.
Counting proper colorings by their nonempty color classes gives
\[
P_G(1)=b_1,\qquad P_G(2)=2b_1+2b_2,\qquad
P_G(3)=3b_1+6b_2+6b_3.
\]
Substitution proves \eqref{eq:general}. For a connected bipartite graph
with an edge, $P_G(1)=0$ and $P_G(2)=2$.
\end{proof}

The array graph $G_n$ has the cells as vertices and exactly the horizontal
and antidiagonal edges defined above. Each row is a path and adjacent rows
share an antidiagonal edge, so $G_n$ is connected. Column parity is a
bipartition. If $p_n=P_{G_n}(3)$, Lemma~\ref{lem:normalization} gives
\begin{equation}\label{eq:anpn}
a_n=\frac{64^n}{48}(p_n+6)\qquad(n\ge1).
\end{equation}

\subsection{The eight boundary states}
Only columns $2,\ldots,6$ of a row constrain the next row.
A proper three-coloring of this five-vertex path, up to permutation of the
three labels, has one of the following eight shapes:
\[
01010,\ 01012,\ 01020,\ 01021,\ 01201,\ 01202,\ 01210,\ 01212.
\]
Here a shape is obtained by replacing labels in order of first occurrence
by $0,1,2$. Every shape has six labeled realizations.
For a representative $s=(s_1,\ldots,s_5)$ and a next row
$q=(q_1,\ldots,q_6)\in\{0,1,2\}^6$, require
\[
q_j\ne q_{j+1},\qquad s_j\ne q_j\qquad(1\le j\le5).
\]
The next state is the shape of $(q_2,\ldots,q_6)$.
Define $T_{st}$ to be the number of these rows with next state $t$.
This definition sums over just $3^6$ words per representative and gives,
in the displayed shape order,
\begin{equation}\label{eq:matrix}
T=\begin{pmatrix}
5&5&4&4&1&1&3&3\\
3&3&3&3&2&2&3&3\\
4&4&4&4&1&1&3&3\\
2&2&2&2&3&3&2&2\\
1&1&1&1&3&3&1&1\\
2&2&2&2&2&2&3&3\\
3&3&3&3&1&1&3&3\\
4&4&2&2&2&2&3&3
\end{pmatrix}.
\end{equation}
For example, the first entry counts the five words
$101010$, $102020$, $121212$, $201010$, $202020$.
For each first-row boundary realization there are two choices for its first
label. Thus every initial shape has weight $12$ and, with $\one$ the
eight-dimensional column of ones,
\begin{equation}\label{eq:pn}
p_n=12\one^{\mathsf T}T^{n-1}\one\qquad(n\ge1).
\end{equation}
The row extension counts all allowable next rows and uses exactly the
information retained in the state, so \eqref{eq:pn} follows by induction
without a height cutoff.

The finite calculation needed for the recurrence is small. Set
$r=T\one$, $t=T^2\one$, $u=T^3\one$. Their transposes are
\begin{align*}
r^{\mathsf T}&=(26,22,24,18,12,18,20,22),\\
t^{\mathsf T}&=(564,456,516,354,222,366,426,462),\\
u^{\mathsf T}&=(11832,9510,10812,7320,4542,7620,8922,9660).
\end{align*}
Direct multiplication of \eqref{eq:matrix} gives
\[
d:=u-25t+90r-81\one=9(-1,1,-1,1,-1,1,-1,1)^{\mathsf T},
\qquad Td=0,\quad\one^{\mathsf T}d=0.
\]
Hence $\one^{\mathsf T}T^m d=0$ for every $m\ge0$, including $m=0$.
By \eqref{eq:pn},
\[
p_{n+3}=25p_{n+2}-90p_{n+1}+81p_n\qquad(n\ge1),
\]
with $(p_1,p_2,p_3)=(96,1944,40392)$. In particular,
\begin{equation}\label{eq:pgf}
P(z):=\sum_{n\ge1}p_nz^n=
\frac{96z-456z^2+432z^3}{1-25z+90z^2-81z^3}.
\end{equation}
Finally \eqref{eq:anpn} yields
\[
\sum_{n\ge1}a_nx^n=\frac{P(64x)}{48}
+\frac{8x}{1-64x},
\]
which simplifies to \eqref{eq:gf}. Multiplying by its denominator gives
\eqref{eq:rec} for $n\ge5$ and the four stated initial values.
The recurrence does not hold at $n=4$ with $a_0=0$.
This proves Theorem~\ref{thm:main}.

\section{Neighboring formulas}
Let $a_{n,w}$ count the same arrays with width $w$, and let $p_{n,w}$
be the corresponding three-color count. For $w\ge2$ the graph is connected
and bipartite, so
\begin{equation}\label{eq:width}
a_{n,w}=2^{nw}(p_{n,w}+6)/48.
\end{equation}
The same transfer construction uses boundaries of length $w-1$.
For $w=2,3,4,5$, the matrices in increasing canonical-shape order are
\[
T_2=(4),\quad T_3=(6),\quad
T_4=\begin{pmatrix}5&5\\4&4\end{pmatrix},\quad
T_5=\begin{pmatrix}5&5&3&3\\4&4&3&3\\2&2&3&3\\4&4&3&3\end{pmatrix}.
\]
The initial weight is $6$ for the single state at width two, and $12$
per state at the other widths. The column sums of $T_4$ are both $9$;
also $T_5^2\one-15T_5\one+18\one=0$.
Consequently
\[
p_{n,2}=6\cdot4^{n-1},\quad
p_{n,3}=12\cdot6^{n-1},\quad
p_{n,4}=24\cdot9^{n-1},
\]
and
\[
\sum_{n\ge1}p_{n,5}z^n=\frac{48z-72z^2}{1-15z+18z^2}.
\]
These explicit matrix identities, together with \eqref{eq:width}, prove
the following corollary, including the first-occurrence normalization.

\begin{corollary}
Barker's formulas in OEIS A233251, A233252 and A233253
\cite{three,four,five} hold for every $n\ge1$. Explicitly,
\begin{align*}
a_{n,3}&=8^{n-1}(6^n+3)/3,&
\sum_{n\ge1}a_{n,3}x^n&=\frac{x(3-64x)}{(1-8x)(1-48x)},\\
a_{n,4}&=2^{4n-3}(4\cdot9^n+9)/9,&
\sum_{n\ge1}a_{n,4}x^n&=\frac{2x(5-208x)}{(1-16x)(1-144x)},\\
&&\sum_{n\ge1}a_{n,5}x^n&=
\frac{4x(9-1120x+30720x^2)}{(1-32x)(1-480x+18432x^2)}.
\end{align*}
Their homogeneous recurrences begin at $n=3,3,4$, respectively.
\end{corollary}

For completeness, column one of the table in A233256 \cite{table}
has an edgeless graph. Lemma~\ref{lem:normalization} gives
\[
a_{n,1}=(6^n+3\cdot4^n+15\cdot2^n)/48\qquad(n\ge1).
\]
Thus the empirical column recurrences for widths $1$ through $6$ in
that entry follow, with coefficient lists
\[
\begin{array}{c|l|c}
w& (c_1,\ldots,c_k)\text{ in }a_{n,w}=\sum_{j=1}^k c_ja_{n-j,w}
&\text{first }n\\\hline
1&(12,-44,48)&4\\
2&(20,-64)&3\\
3&(56,-384)&3\\
4&(160,-2304)&3\\
5&(512,-33792,589824)&4\\
6&(1664,-471040,44826624,-1358954496)&5
\end{array}
\]
The first two row recurrences in A233256 also follow directly. For
width $w\ge2$, a one-row graph is a path and
$p_{1,w}=3\cdot2^{w-1}$. For a two-row graph, classify the last column
by whether its two labels agree. Its transfer is
$\left(\begin{smallmatrix}2&2\\1&1\end{smallmatrix}\right)$,
with initial weights $(3,6)$; its square is three times itself.
Therefore $p_{2,w}=8\cdot3^{w-1}$ for $w\ge2$, and
\[
a_{1,w}=4^w/32+2^w/8,\qquad
a_{2,w}=12^w/18+4^w/8\qquad(w\ge2).
\]
Their recurrences have coefficient lists $(6,-8)$ and $(16,-48)$ and
begin at $w=4$, as specified in the source.

\subsection{The seventh column}
The sixteen shapes of proper six-vertex paths define $T_7$ by exactly the
row-extension rule used above. Every shape has initial weight $12$.
Writing $\one$ for the column of sixteen ones, direct integer multiplication
gives
\[
(T_7^4-42T_7^3+351T_7^2-972T_7+810I)\one=0.
\]
Consequently $p_{n,7}=12\one^{\mathsf T}T_7^{n-1}\one$ satisfies the
corresponding order-four recurrence from $n=5$.
Applying \eqref{eq:width} gives the following additional source result.
\begin{corollary}\label{cor:seven}
Hardin's empirical recurrence in A233255, revision 9, FORMULA \cite{seven},
and the order-five column-seven recurrence in A233256 hold for $n\ge6$:
\begin{align*}
a_{n,7}={}&5504a_{n-1,7}-6438912a_{n-2,7}+2774532096a_{n-3,7}\\
&-478351982592a_{n-4,7}+27831388078080a_{n-5,7}.
\end{align*}
The initial five values are $528$, $1992704$, $8050180096$,
$32835998056448$, and $134179300927602688$.
\end{corollary}
The finite identity specifies its whole matrix: enumerate the proper
words in $\{0,1,2\}^7$, retain their canonical six-letter suffixes, and
count compatible next rows as in \eqref{eq:matrix}.
It is checked entry by entry by the accompanying \texttt{verify.py}.

\subsection{Fixed-height strips and the remaining rows}
Fix a height $h\ge2$. A column now has no internal restrictions.
Let $\mathcal S_h$ be all canonical words of length $h$ on at most three
labels, in lexicographic order. Define the integer matrix $S_h$ by
\begin{equation}\label{eq:column-matrix}
(S_h)_{st}=\#\{q\in\{0,1,2\}^h:\ \operatorname{can}(q)=t,
\ q_r\ne s_r\ (1\le r\le h),\quad q_r\ne s_{r+1}\ (r<h)\}.
\end{equation}
The two inequalities are precisely the horizontal and antidiagonal edges
across consecutive columns. Let $\alpha_h(s)$ be $3$ for the constant
shape and $6$ otherwise. Relabeling symmetry and induction on the number
of columns give
\begin{equation}\label{eq:column-count}
p_{h,w}=\alpha_h^{\mathsf T}S_h^{w-1}\one\qquad(w\ge1).
\end{equation}
Although the original first-occurrence convention is row-major,
Lemma~\ref{lem:normalization} depends on the graph, not its vertex order,
so column transfer can be used without changing the normalization.
For example, the states $000,001,010,011,012$ give
\[
S_3=\begin{pmatrix}2&2&2&2&0\\1&1&0&1&1\\1&1&0&0&0\\
1&1&1&1&0\\0&0&0&1&1\end{pmatrix},\qquad
\alpha_3^{\mathsf T}=(3,6,6,6,6).
\]
To specify the finite identities for every remaining row, put
$g_h(z)=\sum_{j=0}^{d_h}c_{h,j}z^j$ with $c_{h,0}=1$, and define
$g_h^*(S)=\sum_{j=0}^{d_h}c_{h,j}S^{d_h-j}$.
The coefficient lists, in increasing degree, are
\begin{align*}
g_2 &: (1,-3),\quad g_3 : (1,-5,2),\quad g_4 : (1,-9,15,-6),\\
g_5 &: (1,-16,64,-76,-17,84,-48,8),\\
g_6 &: (1,-30,283,-1038,573,6705,-21056,27636,-17216,4128).
\end{align*}
For $g_7$, of degree $27$, the list is split into four consecutive blocks:
\begin{center}\small
\begin{tabular}{rl}
\toprule
Degrees & Coefficients $c_{7,j}$ in increasing degree\\\midrule
$0$--$6$ & $1,-55,1066,-8965,19778,225541,-2020719$\\
$7$--$13$ & $6851286,-7381301,-20655503,79401996,-80882723,-65337447,222811374$\\
$14$--$20$ & $-135953338,-120111164,197698384,-43440564,-76879396,52839536,1810608$\\
$21$--$27$ & $-13170096,4166800,433792,-532928,126464,-13056,512$\\
\bottomrule
\end{tabular}
\end{center}
The matrices defined by \eqref{eq:column-matrix} satisfy the exact identities
\begin{equation}\label{eq:annihilators}
S_h^{t_h}g_h^*(S_h)\one=0,\qquad
\begin{array}{c|rrrrrr}
h&2&3&4&5&6&7\\\hline
|\mathcal S_h|&2&5&14&41&122&365\\
t_h&1&2&4&4&6&6
\end{array}.
\end{equation}
These are finite integer identities, not prefix extrapolations.
The supplementary checker constructs all entries using
\eqref{eq:column-matrix}, evaluates each polynomial on $\one$, multiplies
by $S_h$ exactly $t_h$ times, and checks every coordinate is zero.
Thus all finite premises of \eqref{eq:annihilators} are specified without
large matrix tables.

Here is why these finite identities imply the full infinite sequences.
Equation~\eqref{eq:column-count} and multiplication of
\eqref{eq:annihilators} on the left by
$\alpha_h^{\mathsf T}S_h^m$ show that $p_{h,w}$ satisfies the recurrence
with denominator $g_h(z)$ for $w\ge d_h+t_h+1$.
For $w\ge2$ the array graph is connected, so
$a_{h,w}=2^{hw}(p_{h,w}+6)/48$.
The extra constant $6$ contributes a factor $1-z$, except that
$g_5(1)=0$ already. The exceptional first column is edgeless, with
$a_{h,1}=2^h(3^h+3\cdot2^h+15)/48$.
Consequently, if $F_h(x)=\sum_{w\ge1}a_{h,w}x^w$, the following denominators
give recurrences beginning at the indicated widths:
\begin{center}
\begin{tabular}{cccr}
\toprule
Height $h$ & $R_h(z)$, with denominator $R_h(2^h x)$ & Order & First $w$\\\midrule
2 & $(1-z)g_2(z)$ & 2 & 4\\
3 & $(1-z)g_3(z)$ & 3 & 6\\
4 & $(1-z)g_4(z)$ & 4 & 9\\
5 & $g_5(z)$ & 7 & 12\\
6 & $(1-z)g_6(z)$ & 10 & 17\\
7 & $(1-z)g_7(z)$ & 28 & 35\\
\bottomrule
\end{tabular}
\end{center}
Specifically, if $R_h(z)=\sum_{j=0}^k r_jz^j$, then
$a_{h,w}=-\sum_{j=1}^k r_j2^{hj}a_{h,w-j}$ on the stated range.
The first-column correction can affect coefficients only through degree
$k+1$; the indicated starts also respect that bound.

\begin{corollary}\label{cor:rows}
All empirical row recurrences in the FORMULA block of A233256, revision 6,
hold on their stated ranges. In particular, Barker's formulas in A233257,
revision 8 \cite{rowtwo}, and A233258, revision 8 \cite{rowthree},
and Hardin's recurrence in A233259, revision 6 \cite{rowfour}, hold:
\begin{align*}
F_2(x)&=\frac{x(3-38x+88x^2)}{(1-4x)(1-12x)},\\
F_3(x)&=\frac{x(11-392x+3072x^2-8192x^3+16384x^4)}
{(1-8x)(1-40x+128x^2)},\\
a_{4,w}&=160a_{4,w-1}-6144a_{4,w-2}+86016a_{4,w-3}
-393216a_{4,w-4}\quad(w\ge9).
\end{align*}
Also $a_{2,w}=2^{2w-3}(4\cdot3^w+9)/9$ for $w>1$.
\end{corollary}
\begin{proof}
The previous table gives the recurrences, including the orders $7$, $10$,
and $28$ listed without coefficients in the source.
For the two rational functions, \eqref{eq:column-count} gives the initial
terms $(3,10,104)$ and $(11,136,4672,166400,6049792)$ after normalization.
Multiplication by the stated denominators gives the displayed numerators;
the subsequent coefficients vanish by the recurrences.
The closed form is the two-row formula already established above.
\end{proof}

\subsection{The prior one-column result and cross-references}
Write $b_m$ for A233162 \cite{onecolumn}. Britnell and Wildon's
Proposition 8.1 identifies it with their $B^\dagger_{m-1}(3)$, the number
of set partitions of $m-1$ elements into at most three blocks, with an even
number of elements marked in each block. If
$\left\{\begin{smallmatrix}t\\k\end{smallmatrix}\right\}$ is a Stirling
number of the second kind, each block's parity fixes one binary choice, so
\[
b_m=\sum_{k=0}^3 2^{m-1-k}
\left\{\begin{smallmatrix}m-1\\k\end{smallmatrix}\right\}.
\]
The $k=0$ term is interpreted as $1$ when $m=1$.
For $m\ge2$, substitution of the elementary formulas for the first three
Stirling numbers gives $b_m=a_{m-1,1}$ and
\[
b_m=\frac{2^{m-6}(90+9\cdot2^m+2\cdot3^m)}9,\qquad
\sum_{m\ge1}b_mx^m=
\frac{x(1-11x+35x^2-29x^3)}{(1-2x)(1-4x)(1-6x)}.
\]
Thus Barker's two conjectural formulas of 18 February 2018 in A233162
are consequences of the prior 2015 enumeration; its empirical recurrence
holds for $m\ge5$. This supplies the exact shift for column one of A233256,
without claiming a new enumeration.
The other two identifications cross-referenced by A233256 follow from
the previously recorded formulas in A026244 \cite{knownwidthtwo} and
A007582 \cite{knownrowone}: column two is A026244 at index $n-1$, because
$a_{n,2}=16^n/32+4^n/8$; row one is A007582 at index $w-2$, for $w\ge2$,
because $a_{1,w}=4^w/32+2^w/8$.
For the diagonal sequence A233250 the same graph identity applies;
no fixed-coefficient diagonal recurrence is asserted in that entry.

\section{Coverage and verification}
\begin{center}\small
\begin{tabular}{p{0.32\linewidth}p{0.59\linewidth}}
\toprule
Source item & Mathematical coverage\\\midrule
A233254, FORMULA, revision 8 & Generating function and recurrence proved for all admissible heights.\\
A233251, FORMULA, revision 9 & Both of Barker's formulas and the recurrence proved.\\
A233252, FORMULA, revision 8 & Both of Barker's formulas and the recurrence proved.\\
A233253, FORMULA, revision 7 & Generating function and recurrence proved.\\
A233255, FORMULA, revision 9 & Empirical recurrence proved for heights $n\ge6$.\\
A233257, FORMULA, revision 8 & Both conjectural formulas and the recurrence proved.\\
A233258, FORMULA, revision 8 & Empirical generating function and recurrence proved.\\
A233259, FORMULA, revision 6 & Empirical recurrence proved for widths $w\ge9$.\\
A233256, FORMULA, revision 6 & All listed columns 1--7 and rows 1--7 proved; rows 5--7 use the finite identities \eqref{eq:annihilators}.\\
A233162, FORMULA, revision 21 & Reproved; prior proof Britnell--Wildon (2015), Proposition 8.1 and Section 6.3. Barker's later closed form and generating function follow.\\
A233256, cross-references & Prior formulas cited; three shifted-sequence identifications proved.\\
A233250, revision 8 & Diagonal definition covered by graph normalization; no additional conjecture is stated.\\
\bottomrule
\end{tabular}
\end{center}
The arguments above are written proofs. The accompanying \texttt{verify.py}
enumerates the transitions, verifies the displayed integer identities, compares
the six-column coefficients through degree $30$, and independently counts
the literal six-color arrays through height $6$. It checks the general
normalization formula on every simple graph with one to four vertices.
These finite comparisons supplement the all-height proofs.
For the extended coverage, the checker verifies the sixteen-state seventh-column
identity, all six fixed-height vector identities through dimension $365$,
and the resulting normalized coefficients through height $30$ or width $40$.
The transfer construction and annihilation implication are written proofs;
the larger matrix identities are exact finite computations.

The pinned Lean file
\href{https://github.com/the-omega-institute/trureturing/blob/81ea1b50693b21a1d6370aa3e18815afe5f96af2/D5/S3/Combinatorics/TraceAvoidingArrays/HardinA233254TriangleStrip.lean}{\texttt{HardinA233254TriangleStrip.lean}}
contains declarations in the namespace
\texttt{D5.S3.Combinatorics.TraceAvoidingArrays.}\allowbreak
\texttt{HardinA233254TriangleStrip}.
Its \texttt{result} proves the denominator-multiplied identity
\eqref{eq:gf} in $\mathbb Z[[x]]$ for its literal filtered-array count
\texttt{a}, including $a_0=0$.
Its \texttt{finite\_certificate} gives the exact one-row class counts and
weighted moments for its boundary and first-occurrence flags.
The ordered-graph normalization formula and neighboring-width arguments
above are written proofs; they are not assertions about those two
Lean declarations.
The imported
\href{https://github.com/the-omega-institute/trureturing/blob/81ea1b50693b21a1d6370aa3e18815afe5f96af2/D5/S3/Combinatorics/TraceAvoidingArrays/HardinA233251TriangleStrip.lean}{three-column file}
defines \texttt{widthNormal} and \texttt{widthLegal}, the literal
row-major normalization and edge predicates used by \texttt{a}.
The imported
\href{https://github.com/the-omega-institute/trureturing/blob/81ea1b50693b21a1d6370aa3e18815afe5f96af2/D5/S3/Combinatorics/TraceAvoidingArrays/HardinA233175PairStrip.lean}{pair-strip file}
defines \texttt{allowed}, \texttt{good}, and \texttt{finish}, which track
the earlier occurrences of labels $1$ and $2$.

{\small
\begin{thebibliography}{14}
\setlength{\itemsep}{2pt}
\bibitem{six} R. H. Hardin and C. Barker.
OEIS A233254, revision 8, 11 October 2018.
Definition by Hardin, 6 December 2013; empirical generating function
by Barker, 11 October 2018.
\url{https://oeis.org/A233254}.
\bibitem{three} R. H. Hardin and C. Barker.
OEIS A233251, revision 9, 10 October 2018.
Conjectures in the FORMULA block signed by Barker.
\url{https://oeis.org/A233251}.
\bibitem{four} R. H. Hardin and C. Barker.
OEIS A233252, revision 8, 10 October 2018.
Conjectures in the FORMULA block signed by Barker.
\url{https://oeis.org/A233252}.
\bibitem{five} R. H. Hardin and C. Barker.
OEIS A233253, revision 7, 11 October 2018.
Empirical generating function signed by Barker.
\url{https://oeis.org/A233253}.
\bibitem{table} R. H. Hardin.
OEIS A233256, revision 6, 2 June 2025.
\url{https://oeis.org/A233256}.
\bibitem{seven} R. H. Hardin. OEIS A233255, revision 9, 2 September 2024.
\url{https://oeis.org/A233255}.
\bibitem{rowtwo} R. H. Hardin and C. Barker.
OEIS A233257, revision 8, 11 October 2018.
\url{https://oeis.org/A233257}.
\bibitem{rowthree} R. H. Hardin and C. Barker.
OEIS A233258, revision 8, 11 October 2018.
\url{https://oeis.org/A233258}.
\bibitem{rowfour} R. H. Hardin. OEIS A233259, revision 6, 2 March 2023.
\url{https://oeis.org/A233259}.
\bibitem{onecolumn} R. H. Hardin and C. Barker. OEIS A233162,
revision 21, 31 May 2024. Barker's conjectural formulas, 18 February 2018.
\url{https://oeis.org/A233162}.
\bibitem{knownwidthtwo} N. J. A. Sloane and contributors. OEIS A026244,
revision 47, 9 March 2025. Closed-form definition and Paul Barry's
generating function of 11 March 2004. \url{https://oeis.org/A026244}.
\bibitem{knownrowone} S. Plouffe and contributors. OEIS A007582,
revision 146, 6 November 2024. Closed-form definition and formula block.
\url{https://oeis.org/A007582}.
\bibitem{britnell} J. R. Britnell and M. Wildon.
Bell numbers, partition moves and the eigenvalues of the random-to-top
shuffle in Dynkin types A, B and D. arXiv:1507.04803v1, 16 July 2015,
Proposition 8.1, pp. 23--24, and Section 6.3.
\url{https://arxiv.org/abs/1507.04803v1}.
Journal version: \emph{Journal of Combinatorial Theory, Series A}
\textbf{148} (2017), 116--144. doi:10.1016/j.jcta.2016.12.003.
\bibitem{salas} J. Salas and A. D. Sokal.
Transfer Matrices and Partition-Function Zeros for Antiferromagnetic
Potts Models. I. General Theory and Square-Lattice Chromatic Polynomial.
\emph{Journal of Statistical Physics} \textbf{104} (2001), 609--699.
doi:10.1023/A:1010376605067.
\url{https://arxiv.org/abs/cond-mat/0004330v2}.
\end{thebibliography}
}
\medskip\noindent
\textbf{License.} \copyright\ 2026 Haobo Ma and Wenlin Zhang.
Creative Commons Attribution 4.0 International (CC BY 4.0).
\url{https://creativecommons.org/licenses/by/4.0/}.
\end{document}
